\chapter{SIMD frame diffing}
\label{ch:simd}

Now that we have two grids of packed cells --- the front buffer (what
the terminal is showing) and the back buffer (what we want it to show)
--- we have to compute, quickly, where they differ and emit the
minimal sequence of ANSI bytes that transforms one into the other.
This chapter is the most performance-sensitive corner of \maya{}, and
it is where SIMD earns its keep.

\section{Why SIMD?}

A modern terminal frame is something like 200 columns by 50 rows ---
10,000 cells. At 60 frames per second that is 600,000 cells per
second, and almost all of them are unchanged between consecutive
frames. The naive cell-by-cell comparison is:

\begin{lstlisting}
for (int i = 0; i < width * height; ++i) {
    if (front[i] != back[i]) {
        emit_change(i, back[i]);
    }
}
\end{lstlisting}

One \code{cmp} instruction, one branch, one increment, per cell. The
branches are unpredictable (changes are sparse and clustered), so the
CPU's branch predictor mispredicts often. Cache loads dominate.

\textbf{SIMD} (Single Instruction Multiple Data) lets you compare
multiple cells per instruction. A single AVX2 instruction operates on
a 256-bit register, which is \emph{four} 64-bit cells. AVX-512 doubles
that to eight. NEON, on ARM, gives you two 64-bit lanes per
instruction.

\begin{theory}
SIMD instructions live on dedicated execution units alongside the
scalar ALU. They don't run faster --- they run \emph{wider}. The
throughput win comes from doing more work per instruction without
adding more branches.
\end{theory}

\section{The algorithm}

\maya{}'s diff loop has a beautifully simple shape:

\begin{enumerate}[leftmargin=*]
  \item Iterate row by row within the damage rectangle.
  \item Within a row, use SIMD to skip runs of equal cells four
    (AVX2) or eight (AVX-512) at a time.
  \item When a difference is found, find the first changed cell
    inside the SIMD batch using a bit-scan.
  \item Emit the cursor move, the style change (if any), the
    character bytes.
  \item Continue scanning for the run of \emph{changed} cells until
    they become equal again; then resume bulk-skip.
\end{enumerate}

The hot inner loop is:

\begin{lstlisting}
[[nodiscard]] static std::size_t find_first_diff(
    const uint64_t* a, const uint64_t* b, std::size_t count) noexcept
{
    std::size_t i = 0;
    for (; i + 4 <= count; i += 4) {
        __m256i va  = _mm256_loadu_si256((const __m256i*)(a + i));
        __m256i vb  = _mm256_loadu_si256((const __m256i*)(b + i));
        __m256i cmp = _mm256_cmpeq_epi64(va, vb);
        unsigned mask = _mm256_movemask_pd(_mm256_castsi256_pd(cmp));
        if (mask != 0xFu) {
            return i + std::countr_zero((unsigned)(~mask & 0xFu));
        }
    }
    for (; i < count; ++i) if (a[i] != b[i]) return i;   // scalar tail
    return count;
}
\end{lstlisting}

Read it piece by piece.

\subsection{Loading four cells}

\code{\_mm256\_loadu\_si256} loads 256 bits (four 64-bit cells) from
memory into a register. The \code{u} suffix means ``unaligned'' --- it
tolerates any alignment, with a small performance cost compared to
the aligned variant. \maya{}'s buffer is 64-byte aligned, so in
practice we always hit the fast path.

\subsection{Comparing them}

\code{\_mm256\_cmpeq\_epi64} performs four 64-bit integer
equality checks in parallel. The result is a register where each
lane is \code{0xFFFFFFFFFFFFFFFF} (all ones) if the lanes were equal,
\code{0} otherwise.

\subsection{Reducing to a mask}

\code{\_mm256\_movemask\_pd} extracts the top bit of each 64-bit
lane into a 4-bit integer. (We cast to \code{pd} (packed double)
because \code{movemask} on 64-bit lanes uses the floating-point
variant of the instruction --- same idea, different mnemonic.) Now
\code{mask} is a 4-bit number where bit $k$ is set if lane $k$ was
equal.

\subsection{Detecting any change}

If all four cells were equal, \code{mask == 0xF}. If even one
differs, \code{mask != 0xF}, and we have to find which.

\subsection{Finding the first changed lane}

\code{\textasciitilde mask \& 0xF} flips the bits so that ``not
equal'' is now 1. \code{std::countr\_zero} (count trailing zeros)
returns the index of the lowest set bit. That is the first changed
cell within the batch.

\subsection{The scalar tail}

The SIMD loop steps by 4. If \code{count} is not a multiple of 4, the
last few cells are checked one at a time. This tail is essential ---
forgetting it is the classic SIMD bug.

\section{Architecture dispatch}

Different CPUs support different SIMD widths. \maya{} provides
implementations for:

\begin{itemize}[leftmargin=*]
  \item \textbf{AVX-512.} 8 cells per iteration. Modern Intel server
    chips, recent Xeon, some Ryzen.
  \item \textbf{AVX2.} 4 cells. Haswell (2013) and newer Intel,
    Zen 1+ AMD --- basically every desktop and laptop in current
    use.
  \item \textbf{SSE2.} 2 cells. Universal x86-64 baseline.
  \item \textbf{NEON.} 2 cells per iteration. ARM (Apple Silicon,
    Cortex-A series, AArch64 servers).
  \item \textbf{Scalar.} 1 cell. The portable fallback.
\end{itemize}

The dispatch is done at compile time via the \code{maya::simd::Ops}
template, specialised per ISA tag:

\begin{lstlisting}
struct Avx2 {};
struct Sse2 {};
struct Neon {};
struct Scalar {};

template <typename ISA>
struct Ops {
    static std::size_t find_first_diff(const uint64_t*, const uint64_t*,
                                        std::size_t) noexcept;
    static void streaming_fill(uint64_t*, uint64_t, std::size_t) noexcept;
    // ...
};

template <> struct Ops<Avx2> { /* AVX2 implementation */ };
template <> struct Ops<Sse2> { /* SSE2 implementation */ };
// ...
\end{lstlisting}

The renderer picks the highest-supported ISA detected at build time
(via \code{\#if defined(\_\_AVX2\_\_)} etc.) and uses it everywhere.
At runtime, \maya{} also offers an optional dispatch that selects the
best ISA \emph{on the current CPU} via \code{cpuid} probing --- so a
binary compiled with AVX2 baseline can still light up AVX-512 paths
if the user runs it on a chip that has them.

\begin{idea}
The portable-with-fallback pattern: write the operation once for each
SIMD width, expose the same interface, pick the best at build or
runtime. The architecture-specific code is small and centralised; the
rest of the codebase just calls \code{Ops<ISA>::find\_first\_diff}.
\end{idea}

\section{The other SIMD operations}

\code{find\_first\_diff} is the showpiece, but \maya{}'s SIMD module
has a few more:

\begin{itemize}[leftmargin=*]
  \item \code{skip\_equal(a, b, start, end)} --- bulk-skip equal
    cells inside a row. Same shape as \code{find\_first\_diff} but
    starts at an offset.
  \item \code{streaming\_fill(dst, value, count)} --- fill a region
    with a constant cell (e.g. for clearing). On AVX-512 uses
    non-temporal stores that bypass the cache --- useful for clears,
    because the data won't be read again before being overwritten.
  \item \code{bulk\_eq(a, b, count)} --- whole-row equality, returns
    a bool. Used to check whether an entire row is unchanged, in
    which case the row is skipped in one go.
\end{itemize}

\section{Putting it in the renderer}

The main render loop walks each row of the damage rectangle:

\begin{lstlisting}
for (int y = damage.top; y < damage.bottom; ++y) {
    const uint64_t* front_row = front + y * width;
    const uint64_t* back_row  = back  + y * width;

    std::size_t x = damage.left;
    while (x < damage.right) {
        // Skip equal cells with SIMD.
        x = Ops<ISA>::skip_equal(front_row, back_row, x, damage.right);
        if (x >= damage.right) break;

        // Emit ANSI for the run of changed cells starting at x.
        std::size_t run_end = find_run_end(front_row, back_row, x, damage.right);
        emit_run(out, x, y, back_row + x, run_end - x);
        x = run_end;
    }
}
\end{lstlisting}

The \code{emit\_run} call is where ANSI bytes are written: cursor
positioning, style change (consulting the StylePool's pre-cached
SGR), UTF-8 encoding of the code points. Chapter~\ref{ch:ansi} covers
that side.

\section{What about diff algorithms (Myers, etc.)?}

In version-control land, ``diff'' usually means a clever
longest-common-subsequence algorithm. \maya{}'s diff is not that. It
is line-by-line (or rather cell-by-cell) literal comparison. The
reason: terminal rendering does not benefit from finding ``moved''
content --- moving a row in the terminal still requires writing
the new bytes, because the terminal has no concept of moving a row.
The wall-clock-fastest diff is the one that finds the changed cells
in the most cache-friendly order, which is exactly what packed cells
plus SIMD give us.

\section{Validation}

A subtle bug in any of this would be hard to spot --- the screen
might show a smear of stale glyphs or a corrupted style for a frame.
\maya{} validates the diff in two ways:

\begin{itemize}[leftmargin=*]
  \item Unit tests in \file{maya/tests/test\_diff.cpp} compare
    SIMD output to a known-correct scalar implementation on
    thousands of random inputs.
  \item In debug builds, the canvas can be configured to maintain a
    ``witness chain'' --- a running hash of the bytes emitted, which
    is later replayed in a virtual terminal to confirm that the
    final state matches the back buffer. We will see this in
    Chapter~\ref{ch:inline}.
\end{itemize}

\begin{recap}
The frame differ is \maya{}'s performance heart. It compares the
front and back buffers via SIMD --- four cells per AVX2 instruction,
eight per AVX-512 --- bulk-skips equal runs, and emits ANSI only for
the cells that changed. The pattern (compare, movemask, ctz, scalar
tail) is the canonical SIMD idiom for ``find first $X$ in a stream
of fixed-size records.''
\end{recap}

\section*{Workshop}

\begin{enumerate}[leftmargin=*]
  \item Build and run \file{examples/fps.cpp}. It does a 3D
    raycaster in the terminal. The renderer pushes a fresh frame
    every $\sim$16 ms. Watch your CPU usage --- it should be low.
  \item Read \file{maya/include/maya/platform/simd/avx2.hpp}. Find
    the loop that uses \code{\_mm256\_cmpeq\_epi64}. Identify each
    of the five steps in this chapter.
  \item Write the same \code{find\_first\_diff} in scalar code and
    benchmark both on a million 64-bit integers (random,
    occasionally differing). On a recent x86, the SIMD path should
    be 3-5x faster.
\end{enumerate}
